% ============================================================================
\chapter{System overview}
\label{ch:overview}
% ============================================================================

\section{Hardware components}
\label{sec:system-schema}

The system consists of four hardware elements connected in a chain: the operator's
touch screen, the Raspberry~Pi running this application, the Delta~X~2 controller
board, and the mechanics it drives.

\begin{center}
\begin{tikzpicture}[
  font=\sffamily\small,
  node distance=1.3cm,
  sw/.style={rectangle, rounded corners=2pt, draw=accent, fill=accentlight, thick,
             minimum height=1.5cm, text width=2.7cm, align=center},
  hw/.style={rectangle, rounded corners=2pt, draw=black!45, fill=codegray, thick,
             minimum height=1.5cm, text width=2.7cm, align=center},
  link/.style={{Stealth}-{Stealth}, thick},
]
  \node[hw]                    (screen) {7\,'' touch screen\\[1pt] {\scriptsize 800$\times$480}};
  \node[sw, right=of screen]   (pi)     {Raspberry Pi 5\\[1pt] {\scriptsize \appname{} app}};
  \node[hw, right=of pi]       (ctrl)   {Delta X 2\\ controller};
  \node[hw, right=of ctrl]     (arm)    {delta arms,\\ end effector};
  \draw[link] (screen) -- node[above]{\scriptsize display}
                          node[below]{\scriptsize touch} (pi);
  \draw[link] (pi)     -- node[above]{\scriptsize USB serial}
                          node[below]{\scriptsize G-code} (ctrl);
  \draw[link] (ctrl)   -- node[above]{\scriptsize motor}
                          node[below]{\scriptsize drivers} (arm);
\end{tikzpicture}
\end{center}

\begin{center}
\begin{tabularx}{\textwidth}{@{}lX@{}}
  \toprule
  \textbf{Component} & \textbf{Role} \\
  \midrule
  Delta~X~2 robot & Delta-arm manipulator performing the seeding motion.
    Controlled over a USB serial link (G-code, 115\,200~baud,
    Appendix~\ref{app:gcode}). \\
  Raspberry~Pi~5 & Runs the \appname{} application (Chapter~\ref{ch:rpi}). \\
  7\,'' touch screen & Operator interface, 800\,$\times$\,480 pixels, touch only. \\
  Seeding plates & Interchangeable paper-pot trays; their geometry is described in
    \code{config.toml} (Section~\ref{sec:plates}). \\
  \bottomrule
\end{tabularx}
\end{center}

\section{Software architecture}
\label{sec:architecture}

The application consists of two threads:

\begin{itemize}[itemsep=2pt]
  \item The \textbf{UI thread} runs the Slint event loop and reacts to touch input.
        It never talks to the serial port directly, so the interface stays responsive
        at all times.
  \item The \textbf{robot worker thread} owns the serial connection. It executes one
        command at a time (connect, jog, home, seed a plate) and reports status,
        position and progress back to the UI.
\end{itemize}

\begin{center}
\begin{tikzpicture}[
  font=\sffamily\small,
  comp/.style={rectangle, rounded corners=2pt, draw=accent!75, fill=white, thick,
               align=center, minimum height=0.95cm, text width=3.6cm},
  thread/.style={rectangle, rounded corners=3pt, draw=accent, thick, fill=accentlight,
                 inner sep=10pt},
  shared/.style={rectangle, rounded corners=2pt, draw=warncol, fill=warnlight, thick,
                 align=center, minimum height=0.95cm, text width=3.6cm},
  hw/.style={rectangle, rounded corners=2pt, draw=black!45, fill=codegray, thick,
             align=center, minimum height=0.95cm, text width=3.6cm},
  flow/.style={-{Stealth}, thick},
  bi/.style={{Stealth}-{Stealth}, thick},
]
  % --- Components ---
  \node[comp]                       (loop)    {Slint event loop};
  \node[comp, below=0.45cm of loop] (screens) {screens \& callbacks};
  \node[comp, right=4.9cm of loop]     (robot)  {\code{DeltaRobot}\\[-1pt]
        {\scriptsize G-code, safety limits, position}};
  \node[comp, below=0.45cm of robot]   (serial) {\code{SerialCommunication}};

  % --- Thread boxes (drawn behind the components) ---
  \begin{scope}[on background layer]
    \node[thread, fit=(loop)(screens),
          label={[accent, font=\sffamily\bfseries]above:UI thread}] (uibox) {};
    \node[thread, fit=(robot)(serial),
          label={[accent, font=\sffamily\bfseries]above:robot worker thread}] (wbox) {};
  \end{scope}

  % --- Command / feedback channels ---
  \draw[flow] ([yshift=10pt]uibox.east) -- node[above]{\scriptsize mpsc channel:
        \code{RobotCommand}} ([yshift=10pt]wbox.west);
  \draw[flow] ([yshift=-10pt]wbox.west) -- node[below]{\scriptsize status, position,
        progress} ([yshift=-10pt]uibox.east);

  % --- Shared seeding control flags ---
  \node[shared] (flags) at ($(uibox.south)!0.5!(wbox.south) + (0,-1.5)$)
        {\code{SeedingControl}\\[-1pt] {\scriptsize atomic \code{abort} / \code{pause} flags}};
  \draw[flow] (uibox.south) |- node[pos=0.75, above]{\scriptsize written by
        \textbf{Stop}/\textbf{Pause}} (flags.west);
  \draw[flow] (flags.east) -| node[pos=0.25, above]{\scriptsize read between pots}
        (wbox.south);

  % --- Hardware ---
  \node[hw, below=1.5cm of wbox] (fw) {Delta X 2 firmware};
  \draw[bi] (wbox.east) -- ++(0.5,0) |- node[pos=0.25, right]{\scriptsize USB serial}
        (fw.east);
\end{tikzpicture}
\end{center}

Commands travel from the UI to the worker through a queue and are executed strictly
one at a time. The pause/abort flags take a separate path on purpose: during a
seeding job, the worker is busy inside the job loop and would not look at its queue,
so the \textbf{Stop} and \textbf{Pause} buttons write shared flags that the loop
checks between pots. The robot therefore finishes the pot it is currently working on
and then stops or waits.

\section{Screen map}
\label{sec:screenmap}

\begin{center}
\begin{tikzpicture}[
  font=\sffamily\small,
  scr/.style={rectangle, rounded corners=2pt, draw=accent, fill=accentlight, thick,
              minimum height=0.9cm, text width=2.3cm, align=center},
  nav/.style={-{Stealth}, thick},
  back/.style={-{Stealth}, thick, dashed, black!55},
]
  \node[scr]                                  (main)   {\textbf{Main}};
  \node[scr, above=0.9cm of main]             (calib)  {Calibration};
  \node[scr, below=0.9cm of main]             (system) {System};
  \node[scr, right=1.55cm of main]            (plate)  {Confirm plate};
  \node[scr, right=1.55cm of plate]           (seed)   {Confirm seed};
  \node[scr, right=1.55cm of seed]            (run)    {Seeding};

  \draw[nav]  ([xshift=-8pt]main.north)  -- node[left]{\scriptsize Calibration}
              ([xshift=-8pt]calib.south);
  \draw[back] ([xshift=8pt]calib.south)  -- node[right]{\scriptsize Ok}
              ([xshift=8pt]main.north);
  \draw[nav]  ([xshift=-8pt]main.south)  -- node[left]{\scriptsize Configuration}
              ([xshift=-8pt]system.north);
  \draw[back] ([xshift=8pt]system.north) -- node[right]{\scriptsize Ok}
              ([xshift=8pt]main.south);

  \draw[nav] (main)  -- node[above]{\scriptsize Seeding}       (plate);
  \draw[nav] (plate) -- node[above]{\scriptsize plate}
                        node[below]{\scriptsize chosen}        (seed);
  \draw[nav] (seed)  -- node[above]{\scriptsize Ok}
                        node[below]{\scriptsize (job starts)}  (run);

  \draw[back] (plate.south) to[bend left=18]
              node[pos=0.6, below]{\scriptsize Cancel} (main.south east);
  \draw[back] (seed.south) to[bend left=30]
              node[pos=0.2, below]{\scriptsize Cancel} (main.south east);
  \draw[back] (run.north) to[bend right=20]
              node[above]{\scriptsize Stop / complete / error} (main.north east);
\end{tikzpicture}
\end{center}

\begin{center}
\begin{tabularx}{\textwidth}{@{}lX@{}}
  \toprule
  \textbf{Screen} & \textbf{Purpose} \\
  \midrule
  \textbf{Main} & Home screen; entry to the three functions below. \\
  \textbf{Seeding flow} & Three steps: select the plate type, confirm seeds are
    loaded, then watch/control the running job. \\
  \textbf{Calibration} & Manual jogging, step-size selection, homing, position readout. \\
  \textbf{System} & System utilities (serial-port diagnostics). \\
  \bottomrule
\end{tabularx}
\end{center}

A status bar is visible at the bottom of every screen. It shows a
\textcolor{accent}{$\bullet$}~green\,/\,\textcolor{warncol}{$\bullet$}~red connection
indicator and the most recent status or error message.
